% =========================================================================
% Metaheuristics Project (31-MM34, SoSe 26) -- Full Report (V1)
% Structure skeleton: Report_Structure_V1.tex
% All numbers taken from the repository artifacts:
%   benchmark_logs/*.log, gridsearch_top3_setups.json, ablation.txt,
%   ablation_summary.csv, skillvrp_*_best_solution.csv,
%   skillvrp_*_operator_weights.csv
% Compile: pdflatex Report_Full_V1.tex
% =========================================================================
\documentclass[10pt, a4paper]{article}

\usepackage[utf8]{inputenc}
\usepackage[T1]{fontenc}
\usepackage{lmodern}
\usepackage[margin=2.5cm]{geometry}
\usepackage{amsmath,amssymb}
\usepackage{booktabs}
\usepackage[hidelinks]{hyperref}

\newcommand{\ub}{\%\,UB}

\title{An Adaptive Large Neighborhood Search for the\\
       Prize-Collecting Skill Vehicle Routing Problem}
\author{Umais Chaudhary, Rehman Kerimov, Dustin Schulz, Marco Leander vom Orde \\
        \small 31-MM34 Metaheuristics, SoSe 26, Bielefeld University}
\date{July 2026}

\begin{document}
\maketitle

\begin{abstract}
We develop a metaheuristic for the prize-collecting Skill Vehicle Routing
Problem, in which a team of couriers with heterogeneous skill sets serves an
optional subset of customers under time windows and shift limits so as to
maximize collected profit. The report documents not only the final solver but
the path to it. Guided by the literature on team orienteering and technician
routing, we adopt an Adaptive Large Neighborhood Search (ALNS); as a
yardstick, we first tune an independent single-trajectory LNS to 63.0\,\% of
a singleton upper bound within the 30-second CPython budget. The ALNS
combines 14 destroy and 9 repair operators with penalized simulated-annealing
acceptance, wall-clock-based cooling and runtime-normalized adaptive weights;
its configuration is fixed by a parameter grid search, and a leave-one-out
ablation shows that every operator contributes positively on average. At
equal budget the ALNS matches or beats the tuned baseline on all tuning
instances; ten-minute runs reach 65.5\,\% of the upper bound on average over
all 20 benchmark instances, with every run feasible.
\end{abstract}

% =========================================================================
\section{Introduction}\label{sec:intro}

The pharmacy BiDrug wants to expand its service for elderly customers with
home deliveries, carried out by a small team of couriers. Deliveries differ
in handling requirements (refrigerated transport, controlled substances,
setup of medical devices), and each courier is certified for only a subset of
these \emph{skills}. Because the team is small, not every request can be
served: each served request yields a profit, and the task is to jointly
select and route a subset of customers so that total profit is maximized
while skill compatibility, customer time windows and courier shifts are
respected. The project imposes a hard wall-clock budget (30\,s) on pure
CPython code without third-party packages, and optimal values are unknown, so
solution quality must be judged against a computable upper bound.

Three properties shape the design space. First, the problem is
\emph{selective}: which customers to serve is itself a decision, which
weakens methods that assume all customers are present. Second, skills act as
hard compatibility filters between customers and couriers, making courier
time a heterogeneous, contested resource. Third, the interpreted-Python
budget means that the effective search effort is bounded by how cheaply a
candidate move can be evaluated.

Rather than presenting the final algorithm as a fait accompli, this report
follows the development path, which is also its red thread:
\emph{evidence, yardstick, engineering, breadth, measurement}. The
literature on the two closest problem families points clearly to
destroy-and-repair search (Section~\ref{sec:related}); before building the
full framework we therefore implemented and tuned a deliberately simple
single-trajectory LNS as an independent yardstick, so that every subsequent
design decision could be measured rather than argued
(Section~\ref{sec:baseline-dev}). On top of an evaluation layer engineered
for throughput --- constant-time insertion feasibility and move caching
(Section~\ref{sec:eval}) --- we then added a deliberately broad portfolio of
23 operators and let the data decide: a parameter grid search fixes the
configuration, a leave-one-out ablation quantifies each operator's
contribution, and weight tracking verifies that the adaptive layer allocates
effort sensibly (Section~\ref{sec:components}).

Our contributions are: (i) an ALNS for the prize-collecting Skill VRP with a
skill-aware operator portfolio, including problem-specific skill-scarcity
destroy and scarce-skill-first repair operators; (ii) an evaluation
architecture (forward time slacks, per-route move caching, penalized
scoring) that sustains tens of thousands of destroy--repair iterations per
run in interpreted Python; (iii) a systematic empirical study --- grid
search, 23-operator ablation, operator-weight tracking --- showing that
every operator contributes and that adaptivity prices operators correctly;
and (iv) a transparent evaluation against a tuned LNS baseline at equal
budget and in ten-minute long runs, reaching 65.5\,\ub{} on average with all
runs feasible.

The report follows the four components required by the project description:
the approach is described in Sections~\ref{sec:problem}
and~\ref{sec:method}, the relevant literature is discussed in
Section~\ref{sec:related}, and the experimental investigation of the
components and of the overall performance are the subjects of
Sections~\ref{sec:components} and~\ref{sec:overall}, respectively.

% =========================================================================
\section{Problem Statement}\label{sec:problem}

Let $C=\{1,\dots,n\}$ be the customers and $B$ the couriers; node $0$ is the
pharmacy. Each customer $c$ has profit $p_c$, service time $s_c$, time window
$[e_c,\ell_c]$ and required skills $R_c\subseteq\{1,\dots,S\}$. Each courier
$b$ has a shift $[\mathit{start}_b,\mathit{end}_b]$ and a skill set
$Q_b\subseteq\{1,\dots,S\}$. Travel times $d_{ij}$ are Euclidean distances
rounded to the nearest integer, identical for all couriers. A plan assigns to
every courier a tour that starts and ends at the pharmacy within its shift;
courier $b$ may serve $c$ only if $R_c\subseteq Q_b$; service at $c$ must
begin no later than $\ell_c$ (early arrivals wait until $e_c$); each customer
is visited at most once. With $x_c\in\{0,1\}$ indicating service, the
objective is
\begin{equation}
  \max \sum_{c\in C} p_c\, x_c .
  \label{eq:objective}
\end{equation}
The problem contains the team orienteering problem and hence the TSP as
special cases and is NP-hard. Since optima are unknown, we report quality
relative to the \emph{singleton upper bound}
\begin{equation*}
\mathrm{UB}=\sum_{c\in C_{\mathrm{solo}}} p_c,
\qquad
C_{\mathrm{solo}}=\{\,c\in C : \exists\, b \text{ that can serve } c
\text{ as a one-customer tour}\,\},
\end{equation*}
i.e., the total profit of all customers that at least one courier could serve
alone. UB ignores all interaction between customers and is
therefore loose; it is, however, independent of any solver and thus suitable
for relative comparisons across instances and algorithm versions.

% =========================================================================
\section{Related Work and Choice of Method}\label{sec:related}

\subsection{Problem classification}
The selective aspect --- optional customers, profit maximization under
routing-time restrictions --- places the problem in the orienteering family
\cite{vansteenwegen2011}; with multiple vehicles and hard time windows its
routing core is the team orienteering problem with time windows (TOPTW)
\cite{labadie2012,gunawan2016}. The skill aspect corresponds to the Skill VRP
of Cappanera et al.~\cite{cappanera2011} and to the technician routing and
scheduling problem (TRSP) \cite{kovacs2012,pillac2013}, which combines
skills, time windows and worker shifts exactly as here.

\subsection{Metaheuristics for selective routing with time windows}
The TOPTW literature is dominated by trajectory-based search: the fast
iterated local search of Vansteenwegen et al.~\cite{vansteenwegen2009}
remains the low-runtime reference; granular VNS \cite{labadie2012}, the
three-component heuristic I3CH \cite{hu2014} and a well-tuned ILS/SA hybrid
\cite{gunawan2017} define the quality frontier. A shared efficiency idea in
these works is constant-time insertion feasibility via forward time slacks
\cite{savelsbergh1992,vansteenwegen2009}, which we adopt.

\subsection{LNS, ALNS and population methods}
Large neighborhood search \cite{shaw1998} and its adaptive variant
\cite{ropke2006,pisinger2007} repeatedly destroy part of a solution (random,
related, worst removal) and repair it by greedy or regret insertion; SISR
\cite{christiaens2020} shows that carefully designed removals keep this
family state of the art. For the TRSP --- the closest published relative of
our problem --- Kovacs et al.~\cite{kovacs2012} report ALNS clearly
outperforming exact approaches beyond small instances; the home-health-care
survey of Ciss\'e et al.~\cite{cisse2017} reaches the same conclusion for
compatibility-constrained routing with optional visits. Hybrid genetic
search \cite{vidal2012,vidal2022} defines the state of the art for classical
VRPs, but its strengths (giant-tour decoding, dense route-improvement
neighborhoods) degrade when customer \emph{selection} is part of the problem
and when time windows, shifts and skill filters render most crossover
offspring infeasible.

\subsection{Synthesis}
Both closest problem families point to the same design: destroy-and-repair
search with insertion-based repair, feasibility filtering for skills, and
annealing-type acceptance --- i.e., ALNS. Under a short interpreted-Python
budget, the cheap-move efficiency of LNS-type methods weighs even more
heavily against population methods. This fixes the method family; the rest
of the report is about realizing its potential and \emph{verifying} each
ingredient empirically.

% =========================================================================
\section{Approach Description}\label{sec:method}

\subsection{Development path and baseline LNS}\label{sec:baseline-dev}
We deliberately split development into two solvers. The first,
\texttt{skillvrp+.py}, is a compact single-trajectory LNS: multi-start
greedy insertion construction, random/worst/related removal, greedy
re-insertion and record-to-record-travel acceptance. It was tuned early and
kept frozen as a \emph{yardstick}; over three tuning rounds its 30-second
benchmark average rose from 61.8\,\ub{} to 62.4\,\ub{} and finally
63.0\,\ub{} over all 20 instances (benchmark logs), with every solution
validated by the provided checker. Having a strong, independent reference
this early proved valuable in two ways: it set a non-trivial bar that a more
complex framework had to \emph{earn}, and it exposed that most of the
attainable quality at 30\,s comes from evaluation throughput rather than
from operator sophistication --- the insight that shaped the second solver.

The second solver, \texttt{alns.py}, is the modular ALNS described in the
remainder of this section. It merges the operator ideas of both team
implementations behind common \textsc{Destroy}/\textsc{Repair} interfaces,
so that operators, parameters and instrumentation (Section~\ref{sec:track})
can be exchanged freely.

\subsection{Framework overview}
Algorithm~1 shows the main loop. Operators are drawn with probability
proportional to adaptive weights; candidates are evaluated by a penalized
score and accepted by simulated annealing whose temperature decays with
\emph{elapsed wall-clock time}; weights are updated every segment; after
prolonged stagnation the search restarts from a randomized greedy solution
and reheats. The best feasible solution found is returned; the empty plan is
the always-feasible fallback.

\begin{figure}[t]
\hrule\vspace{3pt}
\noindent\textbf{Algorithm 1:} ALNS for the prize-collecting Skill VRP
\vspace{3pt}\hrule\vspace{4pt}
\small
\begin{tabbing}
\hspace{1.5em}\=\hspace{1.5em}\=\hspace{1.5em}\=\kill
1:\> $x \gets$ best of four greedy insertion orderings;\quad $x^{*} \gets x$\\
2:\> \textbf{while} $\mathrm{now} < \mathrm{deadline}$ \textbf{do}\\
3:\>\> draw destroy $d$ and repair $r$ proportional to weights;\quad
        $x' \gets r(d(x))$\\
4:\>\> $T \gets \max\bigl(T_{\min},\,
        T_0\,(T_{\min}/T_0)^{(\mathrm{now}-t_0)/(\mathrm{deadline}-t_0)}\bigr)$
        \hfill$\triangleright$ time-based cooling\\
5:\>\> \textbf{if} $\mathrm{score}(x') \ge \mathrm{score}(x)$ \textbf{or}
        $\mathrm{rand} < e^{(\mathrm{score}(x')-\mathrm{score}(x))/T}$
        \textbf{then} $x \gets x'$
        \hfill$\triangleright$ penalized score\\
6:\>\> \textbf{if} $x'$ feasible \textbf{and} $f(x') > f(x^{*})$
        \textbf{then} $x^{*} \gets x'$
        \hfill$\triangleright$ optional Or-opt polish hook\\
7:\>\> score the two operators (25/10/3/0);\quad
        every $L$ iterations update weights\\
8:\>\> \textbf{if} 700 consecutive rejections \textbf{then}
        restart from randomized greedy, reheat $T$\\
9:\> \textbf{return} $x^{*}$
\end{tabbing}
\vspace{-6pt}\hrule
\end{figure}

\subsection{Preprocessing and construction heuristic}\label{sec:construction}
Preprocessing computes the rounded distance matrix and, for every customer,
the list of couriers that can serve it \emph{solo} (skill subset test plus a
singleton-tour time check); customers with an empty list are discarded once
and for all. A profit-density bound $p_c/\tau_c$, with $\tau_c$ a lower
bound on the time needed to serve $c$, supports density-based operators.
Construction is greedy insertion at the feasible position of minimal travel
detour over all compatible couriers. Since the insertion order matters, four
orderings are built and evaluated --- by profit, by profit density, by
fewest compatible couriers first, and a lexicographic hybrid --- and the most
profitable start is kept (multi-start greedy).

\subsection{Destroy operators}
Following the taxonomy of common ALNS operators, the portfolio contains 14
destroy operators in five groups. The number of removed customers is sampled
per call from a fraction range of the currently served customers (base range
$U(0.03,0.10)$, related range $U(0.05,0.15)$, capped), so small repairs
polish and large ones restructure.
\begin{itemize}\itemsep2pt
\item \emph{Random}: random removal \cite{ropke2006} (pure
      diversification).
\item \emph{Cost-based}: worst-density removal (static $p_c/\tau_c$
      ranking); two worst-detour variants (lowest profit per travel detour
      in the current position, with biased sampling) \cite{ropke2006}; and
      largest-saving removal (customers whose removal frees the most route
      time).
\item \emph{Relatedness-based}: related removal around a seed customer;
      Shaw removal with a weighted distance/time/skill similarity index
      \cite{shaw1998}; a temporal Shaw variant; and cluster removal, which
      removes one of two MST-based spatial clusters \cite{pisinger2007}.
\item \emph{Structural}: route removal (entire routes), sequence removal
      (contiguous route segments) and time-window-segment removal (all
      served customers in a random time band).
\item \emph{Problem-specific}: skill-scarcity removal targets low-profit
      customers occupying couriers whose skill set is in high demand among
      the currently unserved customers --- freeing contested skill capacity
      that purely geometric operators cannot ``see''. History-based removal
      \cite{ropke2006} removes customers far from their historically best
      position cost.
\end{itemize}

\subsection{Repair operators}
Repair candidates are the removed customers plus a capped sample of
unserved extras (the cap is a tuning parameter,
Section~\ref{sec:exp-grid}); repairing from the full unserved pool was
rejected early because iteration counts collapse. All repairs generate
feasible insertions only. The nine operators are: \emph{greedy best-first}
insertion (globally best insertion, highest profit, ties by minimal
detour); \emph{regret-2} and a courier-level \emph{regret-3} insertion
\cite{ropke2006}, which prioritize customers whose alternatives are much
worse; \emph{sequential cheapest insertion} in profit order
\cite{kovacs2012}, a single-pass variant that trades per-repair care for
iteration throughput; a \emph{noised} greedy variant that perturbs the
position ranking \cite{ropke2006}; \emph{scarce-skill-first} insertion, the
repair counterpart of skill-scarcity removal; \emph{last-removed-first-inserted};
\emph{random-position} insertion; and \emph{Shaw insertion}, which orders
candidates by similarity to the last inserted customer.

\subsection{Efficient move evaluation}\label{sec:eval}
The dominant operation is ``what is the best feasible insertion of customer
$c$ into route $r$?''. Three mechanisms make it cheap enough for
interpreted Python.
\textbf{(i) Route caches with forward time slacks.} Every route caches
arrival, service-start and departure times plus the forward time slack
\cite{savelsbergh1992}: $\mathit{slack}_i$ is the maximal delay of the
service start at position $i$ that keeps the rest of the route and the
depot return feasible, recomputed backwards in $O(m)$ only when the route
changes. An insertion is then checked in $O(1)$: the new customer must meet
its due time and the induced delay of its successor must not exceed the
slack.
\textbf{(ii) Move caching.} Inserting a customer modifies exactly one
route, so cached best moves of all candidates on \emph{other} routes remain
exact. Repairs build a per-(customer, vehicle) cache of top moves once and,
after each insertion, invalidate and recompute only the entries of the
modified vehicle \cite{ropke2006}.
\textbf{(iii) Penalized evaluation.} Candidates are compared by
$\mathrm{score} = \text{profit} - \text{penalties}$, with moderate
per-unit penalties for time-window and shift violations (20 each) and a
prohibitive skill penalty ($10^4$), so the annealing may traverse mildly
infeasible states while skill violations are effectively forbidden; only
feasible solutions can become the incumbent best.
The effect is visible in the iteration counts: at 30\,s the solver performs
roughly 65{,}000 destroy--repair iterations on the 100-customer instance and
still about 2{,}500 on the 1000-customer instance (grid-search logs) ---
enough for the adaptive machinery to act on every instance size.

\subsection{Acceptance, adaptivity and restarts}\label{sec:accept}
Two calibration insights matter more than exact values. First, the initial
temperature must be scaled to the magnitude of a \emph{move}:
destroy--repair deltas are sums of a few customer profits (order $10^2$)
regardless of instance size, so $T_0 = 1.5 \times 100$ in the final
configuration; scaling $T_0$ to the total solution value would turn the
search into a random walk on large instances. Second, cooling is
\emph{time-based} (Algorithm~1, line~4): iteration throughput spans two
orders of magnitude across the benchmark, so any fixed per-iteration rate
is either frozen on small or never cools on large instances. Operator
weights start at 1, are updated every segment (100 iterations) from scores
(new global best 25, improvement 10, accepted 3, rejected 0) smoothed with
reaction factor 0.2, and the scores are \emph{normalized by the operator's
measured runtime} ($\alpha=0.5$), so expensive operators must earn their
cost; a minimum weight keeps every operator alive. After 700 consecutive
rejections the search restarts from a randomized greedy solution and
reheats (factor 0.75). An Or-opt polish-and-refill hook (compress routes,
then greedily refill freed shift capacity) can run on new global bests; the
grid search selected configurations with this hook disabled, indicating
that at these budgets its overhead is not repaid.

% =========================================================================
\section{Experimental Investigation of the Components}\label{sec:components}

\subsection{Setup and evaluation metric}
The benchmark comprises the 20 provided instances with $n=50$ to $1000$
customers, 2--25 couriers, 3--8 skills and a skill-requirement parameter
$k\in\{0,0.5,1,1.5,2\}$ encoded in the instance names. The solver is pure
heuristic Python code --- no exact-solver components --- runs single-threaded
in CPython without third-party packages, conforms to the required command
line (\texttt{python3 skillvrp.py} \emph{instance timeout}) and writes the
tour plan in the specified output format; the empty plan serves as the
always-feasible fallback. Every reported solution passes the provided
checker (\texttt{checker.py}). Quality is reported as profit or as \ub{}
(Section~\ref{sec:problem}). Single 30-second runs scatter by roughly
$\pm 1\,\%$ per instance, so all component studies use at least two seeds
and we treat sub-percent single-seed differences as noise.

\subsection{Initial solution quality}\label{sec:exp-greedy}
The multi-start greedy construction of Section~\ref{sec:construction}
attains on average 49.5\,\ub{} over the 20 instances (range
43.4--55.9\,\ub, ``greedy'' column of Table~\ref{tab:final}), i.e., the
construction alone already secures about half of the loose upper bound with
a guaranteed-feasible solution. Keeping all four insertion orderings is
worthwhile because no single ordering dominates: which ordering yields the
best start varies across instances, and the best-of-four selection is what
the search --- and the restart mechanism, which re-uses randomized greedy
construction --- builds on. The construction's characteristic weakness is
myopic \emph{selection}: it fills shifts in a fixed customer order and
cannot revise which customers deserve the contested courier time. This is
precisely what the destroy--repair search corrects, improving the greedy
start by 20--47\,\% per instance (Section~\ref{sec:overall}); we therefore
invested in search rather than in further construction refinements.

\subsection{Parameter grid search}\label{sec:exp-grid}
The framework's parameters were fixed by a grid search with 30-second runs
on three representative instances (n100, n450, n1000) and two seeds each,
varying the temperature factor, the destroy-size profile, the
unserved-extras cap of the repair candidate set, the penalty profile and the
weight-update parameters. Three observations stand out. First, the
landscape is a \emph{plateau}: the top three configurations differ by less
than 0.3 percentage points in average improvement over the greedy start
(27.4--27.8\,\%), i.e., the method is robust to moderate parameter changes
and the \emph{scale} of $T_0$ matters more than its exact factor. Second,
all top configurations share the \emph{small} destroy profile (base
fraction $U(0.03,0.10)$, related fraction $U(0.05,0.15)$, at most 25\,\% of
served customers): many small, cheap repairs beat few large ones at this
budget. Third, a moderate extras cap (25--40 candidates) suffices; larger
pools cost iterations without paying in quality. The final configuration
uses $T_0=150$, the small destroy profile, cap 40, penalties 20/20/$10^4$,
segment length 100, reaction factor 0.2 and no polish
(Section~\ref{sec:accept}).

\subsection{Operator ablation study}\label{sec:exp-ablation}
Is a 23-operator portfolio bloat? To answer this, a leave-one-out ablation
removed each operator in turn and re-ran the solver on the three
representative instances with two seeds and 90\,s per run (138 runs plus
baseline; segment length 150 and a lower minimum weight were used here so
that weak operators remain visible). Table~\ref{tab:ablation} reports, for
each operator, the average and worst change relative to the all-operator
baseline. The central finding: \emph{removing any operator hurts on
average} --- the mean impact is negative for all 23 operators, so the
portfolio contains no dead weight, and the runtime-normalized adaptive
weights successfully amortize even rarely useful operators. The ranking is
informative: the strongest contributors are the solution-dependent
cost-based and relatedness-based removals (worst-detour~v2 $-1.7\,\%$,
Shaw-related $-1.1\,\%$, largest-saving $-1.0\,\%$), confirming that
coherent, cost-aware destruction is what enables the repairs to restructure
a region. The problem-specific skill-scarcity removal ranks fourth
($-0.8\,\%$, worst case $-1.6\,\%$), validating the design intuition that
skill capacity --- not geometry --- is the contested resource. On the repair
side, greedy best-first insertion is the most valuable single operator
($-0.6\,\%$); the remaining repairs form a long tail of small but
consistently positive contributions.

\begin{table}[t]
\centering\small
\caption{Leave-one-out ablation (3 instances $\times$ 2 seeds $\times$
90\,s). Change in profit vs.\ the all-operator baseline when the operator is
removed; negative = the operator helps.}
\label{tab:ablation}
\begin{tabular}{llrr@{\qquad}llrr}
\toprule
operator & type & avg & worst & operator & type & avg & worst \\
\midrule
worst-detour rem.\ v2 & destr. & $-1.7$ & $-4.7$ & noisy greedy ins.    & rep. & $-0.3$ & $-1.2$ \\
Shaw-related rem.     & destr. & $-1.1$ & $-2.9$ & scarce-skill-first ins. & rep. & $-0.3$ & $-1.3$ \\
largest-saving rem.   & destr. & $-1.0$ & $-2.9$ & Shaw ins.            & rep. & $-0.3$ & $-1.4$ \\
skill-scarcity rem.   & destr. & $-0.8$ & $-1.6$ & regret-3 courier ins.& rep. & $-0.2$ & $-0.4$ \\
worst-detour rem.     & destr. & $-0.7$ & $-1.6$ & random-position ins. & rep. & $-0.2$ & $-1.3$ \\
cluster rem.          & destr. & $-0.7$ & $-2.0$ & random rem.          & destr. & $-0.2$ & $-0.7$ \\
time-segment rem.     & destr. & $-0.7$ & $-2.1$ & related rem.         & destr. & $-0.2$ & $-2.0$ \\
temporal Shaw rem.    & destr. & $-0.7$ & $-1.6$ & LRFI ins.            & rep. & $-0.2$ & $-1.0$ \\
greedy best ins.      & rep.   & $-0.6$ & $-1.7$ & worst-density rem.   & destr. & $-0.2$ & $-1.2$ \\
route rem.            & destr. & $-0.5$ & $-1.3$ & history-based rem.   & destr. & $-0.1$ & $-1.3$ \\
regret-2 ins.         & rep.   & $-0.4$ & $-1.9$ & & & & \\
sequential ins.\ (profit) & rep. & $-0.4$ & $-1.4$ & & & & \\
sequence rem.         & destr. & $-0.4$ & $-1.5$ & & & & \\
\bottomrule
\end{tabular}
\end{table}

\subsection{Adaptive behavior}\label{sec:track}
An instrumented variant of the solver (\texttt{alns\_tracking.py}) records
the best-solution timeline and the full weight trajectory of every operator
per 100-iteration segment. The traces confirm that the adaptive layer works
as intended rather than merely averaging: weights start uniform at 1 and
differentiate within the first segments; on the 1000-customer instance the
final segments are led by worst-detour removal (0.79), courier-level
regret-3 insertion (0.74) and largest-saving removal (0.71) --- broadly
consistent with the ablation ranking obtained by a completely different
protocol --- while the cheap sequential insertion retains a substantial
share through its runtime normalization. Restarts were never triggered in
the tracked benchmark runs, i.e., the SA schedule kept the acceptance rate
above the stagnation threshold for the whole budget.

% =========================================================================
\section{Overall Performance}\label{sec:overall}

\subsection{Comparison with the tuned LNS baseline at equal budget}
At the project budget of 30\,s we compare the final ALNS configuration
against the frozen LNS yardstick of Section~\ref{sec:baseline-dev}
(benchmark average 63.0\,\ub) on the three grid-search instances
(Table~\ref{tab:head2head}). The ALNS matches the baseline on the small
instance and is clearly ahead on the mid-size and large ones ($+1.7\,\%$ and
$+1.9\,\%$) --- evidence that the overhead of 23 operators, adaptive weights
and penalized evaluation is fully amortized by the cached $O(1)$ move
evaluation, and that the framework converts the extra machinery into
quality exactly where the search space is largest.

\begin{table}[t]
\centering\small
\caption{30-second head-to-head: mean ALNS profit over two seeds
(final configuration, grid-search protocol) vs.\ the single-run tuned LNS
baseline.}
\label{tab:head2head}
\begin{tabular}{lrrr}
\toprule
instance & ALNS (mean of 2 seeds) & LNS baseline & difference \\
\midrule
n100\_s5\_k1.0  &  3\,006 &  3\,004 & $+0.1\,\%$ \\
n450\_s3\_k0.0  & 14\,393 & 14\,155 & $+1.7\,\%$ \\
n1000\_s8\_k2.0 & 36\,462 & 35\,781 & $+1.9\,\%$ \\
\bottomrule
\end{tabular}
\end{table}

\subsection{Long-run results}
Table~\ref{tab:final} reports ten-minute runs of the final configuration on
all 20 instances, together with the singleton upper bound, the greedy start
and, for reference, the 30-second LNS baseline. All runs are feasible and
pass the checker. The average quality is \textbf{65.5\,\ub}; the ALNS
improves its greedy start by 20--47\,\% (average 32.5\,\%), underlining that
the search --- not the construction --- carries the result. Compared with
the 30-second baseline (63.0\,\ub{} average), the long-run ALNS is ahead on
19 of 20 instances; the single exception is n150, where the baseline's
4\,936 edges the ALNS's 4\,930. Note that this last comparison spans
different budgets and is reported for orientation, not as a head-to-head
claim; the equal-budget comparison is Table~\ref{tab:head2head}.

\begin{table}[t]
\centering\small
\caption{Long-run results (600\,s, final configuration): singleton upper
bound, greedy start, ALNS best profit, quality, time of last improvement
and iterations. Rightmost column: 30\,s LNS baseline for orientation.}
\label{tab:final}
\begin{tabular}{lrrrrrrrr}
\toprule
instance & $n$ & $B$ & UB & greedy & ALNS & \ub & $t^{*}$[s] & LNS (30\,s) \\
\midrule
n50\_s3\_k0.0   &   50 &  2 &  2\,717 &  1\,501 &  1\,815 & 66.8 &  22 &  1\,793 \\
n75\_s4\_k0.5   &   75 &  2 &  4\,255 &  2\,058 &  2\,798 & 65.8 & 266 &  2\,784 \\
n100\_s5\_k1.0  &  100 &  2 &  5\,160 &  2\,292 &  3\,050 & 59.1 &  27 &  3\,004 \\
n125\_s6\_k1.5  &  125 &  3 &  6\,866 &  3\,697 &  4\,444 & 64.7 & 104 &  4\,223 \\
n150\_s8\_k2.0  &  150 &  3 &  8\,259 &  4\,065 &  4\,930 & 59.7 &  37 &  4\,936 \\
n200\_s3\_k0.0  &  200 &  5 & 10\,236 &  5\,504 &  7\,573 & 74.0 & 319 &  7\,487 \\
n250\_s4\_k0.5  &  250 &  6 & 14\,034 &  7\,845 & 10\,060 & 71.7 & 136 &  9\,874 \\
n300\_s5\_k1.0  &  300 &  7 & 16\,729 &  8\,331 & 11\,770 & 70.4 & 280 & 11\,450 \\
n350\_s6\_k1.5  &  350 &  8 & 19\,050 &  9\,708 & 12\,978 & 68.1 & 318 & 12\,583 \\
n400\_s8\_k2.0  &  400 & 10 & 22\,335 & 10\,426 & 13\,198 & 59.1 & 557 & 12\,741 \\
n450\_s3\_k0.0  &  450 & 11 & 24\,478 & 11\,003 & 14\,890 & 60.8 & 169 & 14\,155 \\
n500\_s4\_k0.5  &  500 & 12 & 27\,820 & 13\,344 & 18\,498 & 66.5 & 450 & 17\,620 \\
n600\_s5\_k1.0  &  600 & 15 & 32\,477 & 16\,277 & 22\,016 & 67.8 & 314 & 21\,177 \\
n700\_s6\_k1.5  &  700 & 17 & 37\,852 & 18\,732 & 25\,140 & 66.4 & 470 & 23\,758 \\
n750\_s8\_k2.0  &  750 & 18 & 40\,329 & 19\,542 & 26\,785 & 66.4 & 551 & 25\,217 \\
n800\_s3\_k0.0  &  800 & 20 & 43\,128 & 20\,995 & 27\,131 & 62.9 & 485 & 25\,491 \\
n850\_s4\_k0.5  &  850 & 21 & 46\,190 & 21\,708 & 31\,796 & 68.8 & 577 & 29\,875 \\
n900\_s5\_k1.0  &  900 & 22 & 49\,703 & 23\,101 & 30\,243 & 60.8 & 587 & 28\,660 \\
n950\_s6\_k1.5  &  950 & 23 & 52\,946 & 22\,960 & 31\,545 & 59.6 & 447 & 29\,349 \\
n1000\_s8\_k2.0 & 1000 & 25 & 54\,664 & 30\,193 & 38\,508 & 70.4 & 591 & 35\,781 \\
\midrule
average &&&&&& \textbf{65.5} && (63.0) \\
\bottomrule
\end{tabular}
\end{table}

\subsection{Convergence and scaling}
The time of the last improvement $t^{*}$ in Table~\ref{tab:final} separates
the benchmark cleanly: small instances ($n\le 150$) essentially converge
within the first 30--40 seconds, while instances with $n\ge 400$ keep
improving until the very end of the ten-minute budget ($t^{*}>440\,$s on 9
of the 10 largest instances) --- large instances remain compute-bound, and
additional budget translates directly into profit. Iteration counts in
600\,s range from about 85{,}000 ($n=150$) to 860{,}000 ($n=75$) and are
still close to 100{,}000 on $n=1000$. Solution quality shows \emph{no}
systematic degradation with instance size --- the best relative results
occur at $n=200$ (74.0\,\ub) and $n=1000$ (70.4\,\ub) --- whereas the
weakest instances (around 59--61\,\ub) combine few couriers with elevated
skill requirements, consistent with skill capacity rather than geometry
being the binding resource. Robustness requirements are met throughout: all
runs, at 30\,s and 600\,s, returned feasible, checker-validated solutions
within the limit, and no run needed the restart fallback.

% =========================================================================
\section{Conclusion}

We approached the prize-collecting Skill VRP by fixing the method family
from the literature, building a tuned single-trajectory LNS as an early
yardstick, engineering the move evaluation for throughput, and only then
adding breadth --- 14 destroy and 9 repair operators whose configuration and
usefulness were established by grid search, a 23-operator leave-one-out
ablation and weight tracking rather than by intuition. The result is a
solver that matches or beats the tuned baseline at the 30-second project
budget, reaches 65.5\,\% of the singleton upper bound on average in
ten-minute runs, and returned a feasible solution in every single run. The
methodological lessons carry beyond this problem: under short interpreted
budgets, evaluation throughput is the enabling resource for adaptivity; a
frozen baseline turns design debates into measurements; and
runtime-normalized adaptive weights make a broad operator portfolio an
asset rather than bloat --- the ablation found no operator whose removal
helps on average. Promising next steps are string-based removals in the
spirit of SISR \cite{christiaens2020}, a time-aware insertion score
$p_c-\lambda\,\Delta\mathrm{travel}$ reflecting that shift time is the
scarce resource, and a set-covering post-optimization over the route pool
collected during the search \cite{hu2014}.

% =========================================================================
\begin{thebibliography}{19}\small

\bibitem{cappanera2011} P.~Cappanera, L.~Gouveia, M.G.~Scutell\`a.
The skill vehicle routing problem. \emph{Network Optimization} (INOC 2011),
LNCS 6701, 354--364, 2011.

\bibitem{christiaens2020} J.~Christiaens, G.~Vanden Berghe.
Slack induction by string removals for vehicle routing problems.
\emph{Transportation Science} 54(2), 417--433, 2020.

\bibitem{cisse2017} M.~Ciss\'e, S.~Yal\c{c}{\i}nda\u{g}, Y.~Kergosien,
E.~\c{S}ahin, C.~Lent\'e, A.~Matta. OR problems related to home health care:
a review of relevant routing and scheduling problems.
\emph{Operations Research for Health Care} 13--14, 1--22, 2017.

\bibitem{gunawan2016} A.~Gunawan, H.C.~Lau, P.~Vansteenwegen.
Orienteering problem: a survey of recent variants, solution approaches and
applications. \emph{European Journal of Operational Research} 255(2),
315--332, 2016.

\bibitem{gunawan2017} A.~Gunawan, H.C.~Lau, P.~Vansteenwegen, K.~Lu.
Well-tuned algorithms for the team orienteering problem with time windows.
\emph{Journal of the Operational Research Society} 68(8), 861--876, 2017.

\bibitem{hu2014} Q.~Hu, A.~Lim. An iterative three-component heuristic for
the team orienteering problem with time windows.
\emph{European Journal of Operational Research} 232(2), 276--286, 2014.

\bibitem{kovacs2012} A.A.~Kovacs, S.N.~Parragh, K.F.~Doerner, R.F.~Hartl.
Adaptive large neighborhood search for service technician routing and
scheduling problems. \emph{Journal of Scheduling} 15(5), 579--600, 2012.

\bibitem{labadie2012} N.~Labadie, R.~Mansini, J.~Melechovsk\'y,
R.~Wolfler Calvo. The team orienteering problem with time windows: an
LP-based granular variable neighborhood search.
\emph{European Journal of Operational Research} 220(1), 15--27, 2012.

\bibitem{pillac2013} V.~Pillac, C.~Gu\'eret, A.L.~Medaglia. A parallel
matheuristic for the technician routing and scheduling problem.
\emph{Optimization Letters} 7(7), 1525--1535, 2013.

\bibitem{pisinger2007} D.~Pisinger, S.~Ropke. A general heuristic for
vehicle routing problems. \emph{Computers \& Operations Research} 34(8),
2403--2435, 2007.

\bibitem{ropke2006} S.~Ropke, D.~Pisinger. An adaptive large neighborhood
search heuristic for the pickup and delivery problem with time windows.
\emph{Transportation Science} 40(4), 455--472, 2006.

\bibitem{savelsbergh1992} M.W.P.~Savelsbergh. The vehicle routing problem
with time windows: minimizing route duration.
\emph{ORSA Journal on Computing} 4(2), 146--154, 1992.

\bibitem{shaw1998} P.~Shaw. Using constraint programming and local search
methods to solve vehicle routing problems. \emph{Principles and Practice of
Constraint Programming} (CP 1998), LNCS 1520, 417--431, 1998.

\bibitem{vansteenwegen2009} P.~Vansteenwegen, W.~Souffriau,
G.~Vanden Berghe, D.~Van Oudheusden. Iterated local search for the team
orienteering problem with time windows.
\emph{Computers \& Operations Research} 36(12), 3281--3290, 2009.

\bibitem{vansteenwegen2011} P.~Vansteenwegen, W.~Souffriau,
D.~Van Oudheusden. The orienteering problem: a survey.
\emph{European Journal of Operational Research} 209(1), 1--10, 2011.

\bibitem{vidal2012} T.~Vidal, T.G.~Crainic, M.~Gendreau, N.~Lahrichi,
W.~Rei. A hybrid genetic algorithm for multidepot and periodic vehicle
routing problems. \emph{Operations Research} 60(3), 611--624, 2012.

\bibitem{vidal2022} T.~Vidal. Hybrid genetic search for the CVRP:
open-source implementation and SWAP* neighborhood.
\emph{Computers \& Operations Research} 140, 105643, 2022.

\end{thebibliography}

\end{document}
